\chapter{Instruction encoding and conventions}\label{ch:encoding}

\section{Instruction formats}\label{sec:formats}

Every instruction is one 32-bit little-endian word, aligned on a 4-byte boundary. Bits 31:26 hold the
primary opcode \field{op}. The remaining bits follow one of six formats.

\newcommand{\fmtrow}[2]{\par\noindent\textbf{#1}\hfill\small #2\par\nobreak}

\fmtrow{R: register}{ALU, TAMO, CAP, CAP2, XMAP, XCALL}
\begin{encoding}
\bitheader{0,10,11,15,16,20,21,25,26,31}\\
\fv{6}{op}\fv{5}{rd}\fv{5}{rs1}\fv{5}{rs2}\fv{11}{funct}
\end{encoding}

\fmtrow{I: immediate}{ALU immediates, LUI, loads, JALR, TWAIT, TPOLL, TANY, CSRR, CSRW, SYS, LTAG, LINST}
\begin{encoding}
\bitheader{0,15,16,20,21,25,26,31}\\
\fv{6}{op}\fv{5}{rd}\fv{5}{rs1}\fv{16}{imm16}
\end{encoding}

\fmtrow{S: store}{SD, SW, SH, SB}
\begin{encoding}
\bitheader{0,15,16,20,21,25,26,31}\\
\fv{6}{op}\fv{5}{rs2 (data)}\fv{5}{rs1 (base)}\fv{16}{imm16}
\end{encoding}

\fmtrow{B: branch}{BEQ, BNE, BLT, BGE, BLTU, BGEU}
\begin{encoding}
\bitheader{0,15,16,20,21,25,26,31}\\
\fv{6}{op}\fv{5}{rs1}\fv{5}{rs2}\fv{16}{off16 (instructions)}
\end{encoding}

\fmtrow{J: jump}{JAL}
\begin{encoding}
\bitheader{0,20,21,25,26,31}\\
\fv{6}{op}\fv{5}{rd}\fv{21}{off21 (instructions)}
\end{encoding}

\fmtrow{X: transfer}{TGET, TPUT, TGETX}
\begin{encoding}
\bitheader{0,4,5,9,10,13,14,18,19,22,23,25,26,31}\\
\fv{6}{op}\fv{3}{sub}\fv{4}{tag}\fv{5}{rl}\fv{4}{cs}\fv{5}{ro}\fv{5}{rn}
\end{encoding}

\subsection{Fields}
\begin{description}[leftmargin=1em,style=nextline]
\item[\field{rd}, \field{rs1}, \field{rs2}] Integer registers, in bits 25:21, 20:16 and 15:11. In the S and
  B formats the register in bits 25:21 is a \emph{source}.
\item[\field{cd}, \field{cs}, \field{ce}] Capability registers (\field{ce} names an endpoint). When a
  capability instruction names a capability in the \field{rd} or \field{rs1} position, the low four bits
  select \reg{c0}--\reg{c15} and the high bit (25 or 20) must be zero, else the instruction is
  \cause{illegal}. Transfers, atomics and overlay instructions have dedicated 4-bit \field{cs} fields.
\item[\field{imm16}] Sign-extended for arithmetic, comparisons, loads, stores and \insn{jalr};
  zero-extended for \insn{andi}, \insn{ori} and \insn{xori}; a 16-bit CSR number for \insn{csrr}/\insn{csrw};
  a tag mask for \insn{twait}/\insn{tpoll}.
\item[\field{off16}, \field{off21}] Signed offsets counted in instructions (multiplied by 4) relative to
  the address of the branch or jump itself.
\item[\field{funct}] Selects the operation within an opcode. For CAP, bits 3:0 select the function, bit 4
  is a flag (sender filter for \insn{ctake}) and bits 10:5 hold the save-table slot. For CAP2, bits 3:0
  select the function and bits 10:4 are ignored. For XMAP and XCALL, bit 4 selects the second instruction
  of the pair (\insn{xunmap}, \insn{xret}) and bits 3:0 hold \field{cs}.
\end{description}

Fields shown shaded in the instruction entries are \term{reserved}. Implementations ignore them: an
instruction executes identically whatever they hold. Software must write them as zero, because future
editions may assign them. Fields shown in colour are fixed opcode or function values, and they are
checked: an encoding whose fixed fields match no entry is \cause{illegal}. Table~\ref{tab:fieldrules}
lists, for every opcode, which bits are checked and which are ignored.

\begin{table}[htbp]
\centering\small
\caption{Checked and ignored encoding bits}\label{tab:fieldrules}
\begin{tabularx}{\linewidth}{@{}l X X@{}}
\toprule
\textbf{Opcode} & \textbf{Checked (else \cause{illegal})} & \textbf{Ignored (reserved)}\\\midrule
ALU & \field{funct}\texttt{<10:4>} = 0 & ---\\
LUI & --- & bits 20:16\\
XFER & \field{sub} $\le$ 2 & ---\\
TWAIT, TPOLL, CSRR, CSRW & --- & the unused register field\\
TANY, LTAG, LINST & --- & \field{imm16}\\
TAMO & \field{funct}\texttt{<2:0>} $\le$ 4, \field{funct}\texttt{<6:3>} = 0 & ---\\
CAP & bit 25 = 0 where bits 24:21 name a capability; bit 20 = 0 where bits 19:16 name a capability &
  unused register fields; \field{funct}\texttt{<10:5>} except as a save-table slot;
  \field{funct}\texttt{<4>} except for \insn{ctake}\\
CAP2 & bits 25 and 20 = 0; \field{funct}\texttt{<3:0>} $\le$ 5 & \field{funct}\texttt{<10:4>}; unused register fields\\
SYS & \field{imm16} $\le$ 2 & bits 25:16\\
XMAP & map index (bits 25:21) $\le$ 15 & unused bits of \field{funct}; for \insn{xunmap}, bits 20:11\\
XCALL & --- & unused bits of \field{funct} and unused register fields\\
\bottomrule
\end{tabularx}
\end{table}

\section{Opcode map}

\input{gen-opmap}

Within an opcode, the function assignments are:

\begin{center}\small
\begin{tabular}{@{}l l l@{}}
\toprule
\textbf{Opcode} & \textbf{Selector} & \textbf{Values}\\\midrule
ALU & \field{funct}\texttt{<3:0>}, bits 10:4 zero & 0 add, 1 sub, 2 and, 3 or, 4 xor, 5 sll, 6 srl, 7 sra,\\
 & & 8 slt, 9 sltu, 10 mul, 11 mulhu, 12 divu, 13 remu, 14 div, 15 rem\\
XFER & \field{sub} (bits 25:23) & 0 tget, 1 tput, 2 tgetx; 3--7 illegal\\
TAMO & \field{funct}\texttt{<2:0>}, bits 6:3 zero & 0 add, 1 swap, 2 cas, 3 and, 4 or; 5--7 illegal\\
CAP & \field{funct}\texttt{<3:0>} & 0 cmov, 1 csplit, 2 cjoin, 3 cshare, 4 catomic, 5 crestrict,\\
 & & 6 cbase, 7 clen, 8 cperms, 9 cgive, 10 ctake, 11 cclear,\\
 & & 12 csave, 13 crestore, 14 cvalid, 15 cpoll\\
CAP2 & \field{funct}\texttt{<3:0>} & 0 cdom, 1 crevoke, 2 cundom, 3 csend, 4 crecv, 5 cmkep; 6--15 illegal\\
SYS & \field{imm16} & 0 halt, 1 eret, 2 svc; all others illegal\\
XMAP & \field{funct}\texttt{<4>} & 0 xmap, 1 xunmap\\
XCALL & \field{funct}\texttt{<4>} & 0 xcall, 1 xret\\
\bottomrule
\end{tabular}
\end{center}

Appendix~\ref{app:index} lists every instruction alphabetically with its opcode, function and page.

\section{Pseudocode conventions}\label{sec:pseudo}

The \emph{Operation} section of each entry is written in a small pseudocode language.

\begin{center}\small
\begin{tabularx}{\linewidth}{@{}l X@{}}
\toprule
\textbf{Notation} & \textbf{Meaning}\\\midrule
\texttt{X[r]} & integer register $r$; reading \texttt{X[0]} gives 0 and writing it has no effect\\
\texttt{C[c]}, \texttt{CTAB[s]} & capability register $c$, save-table slot $s$\\
\texttt{PC}, \texttt{NextPC} & address of this instruction; address of the next one (default \texttt{PC + 4})\\
\texttt{EPC}, \texttt{TVEC}, \ldots & CSRs, by name; the banked trap CSRs name the current mode's copy\\
\texttt{SEPC}, \texttt{UEPC}, \ldots & the supervisor's and the user's copy of a banked trap CSR\\
\texttt{Mode} & the current privilege mode, \texttt{Supervisor} or \texttt{User} (\csr{status} bit 0)\\
\texttt{DOM[d]} & entry $d$ of the machine-wide revocation-domain table (Section~\ref{sec:domains})\\
\texttt{DEV}, \texttt{CORE}, \texttt{EP} & spaces 4 (device), 5 (core control) and 6 (endpoint)\\
\texttt{instr} & the instruction word\\
\texttt{v<h:l>} & bits $h$..$l$ of $v$\\
\texttt{SInt}, \texttt{UInt} & a bit vector as a signed or unsigned integer\\
\texttt{SignExtend(v, n)}, \texttt{ZeroExtend(v, n)} & extend $v$ to $n$ bits\\
\texttt{ASR(v, n)} & arithmetic shift right\\
\texttt{DLS.Read(a, n)}, \texttt{DLS.Write(a, n, v)} & little-endian access to $n$ bytes of the D-LS\\
\texttt{Far.Read(sp, a, n)}, \texttt{Far.Write(\ldots)} & access to far space \texttt{sp} at absolute address \texttt{a}\\
\texttt{Device.Read(port, i)}, \texttt{Device.Write(port, i, b)} & beat $i$ of a device transfer starting at \texttt{port}: one byte from or to device offset \texttt{port} + $i$, or the next byte of the stream at \texttt{port} (Section~\ref{sec:devxfer})\\
\texttt{Trap(cause, value)} & take a trap (Chapter~\ref{ch:traps}); the instruction has no other effect\\
\texttt{Stall} & wait; no architectural effect\\
\texttt{Block} & do not execute now; retry when another node has made progress\\
\texttt{DrainTransfers()} & wait until all of this core's in-flight transfers have completed\\
\texttt{WaitMemory()} & wait until none of this core's transfers to T2 or T3 is in flight (the publication rule)\\
\texttt{WaitPendingReads(ls, n)} & if $n \neq 0$ and D-LS range $[\textit{ls}, \textit{ls}+n)$ overlaps the destination of an in-flight device read of this core: \texttt{Trap(pend, ls)} when CCR.PTRAP is set, otherwise wait until the read completes\\
\texttt{BusyTags()} & 16-bit mask of tags with transfers in flight (in $\mathcal{S}$, only pending device reads)\\
\texttt{N} & number of cores; node \texttt{N} is block device 0\\
\bottomrule
\end{tabularx}
\end{center}

All arithmetic is modulo $2^{64}$ unless written with \texttt{SInt}/\texttt{UInt}. The steps of an
operation that precede a \texttt{Trap} have no lasting effect.

\subsection{Shared functions}

\begin{pseudo}
// A capability can be used only while valid and, if it belongs to a revocation
// domain, while that domain is alive and has not been revoked since it was made.
func Live(c)
    if not c.valid or c.dom == 0 then return c.valid; end
    if c.V then return DOM[c.dom].alive; end        // a revoker outlives revocations
    return DOM[c.dom].alive and DOM[c.dom].gen == c.gen;
end
func IsRevoker(c)   = c.valid and c.V and c.dom != 0 and DOM[c.dom].alive
func Linear(c)      = not c.D                  // linear unless copyable
func Replaceable(c) = not Live(c) or c.D       // may be overwritten

// Issue-time checks of TGET (kind GET, need R), TPUT (PUT, W) and TGETX (GETX, X).
func CheckTransfer(cap, kind, ls, off, n, lslimit)
    if not Live(cap) then Trap(capinv, off); end
    if not cap.need then Trap(perm, off); end
    if cap.space == EP then Trap(perm, off); end
    if cap.space in {DEV, CORE} and kind == GETX then Trap(perm, off); end
    if cap.space == CORE and kind != PUT then Trap(perm, off); end
    if cap.space != DEV and (ls OR off OR n) MOD 8 != 0 then Trap(xalign, off); end
    if cap.space == DEV then
        if UInt(off) >= cap.len then Trap(bounds, off); end      // off names a port
    elsif UInt(off) > cap.len or UInt(n) > cap.len - off then
        Trap(bounds, off);
    end
    if UInt(ls) > lslimit or UInt(n) > lslimit - ls then Trap(lsfault, ls); end
end

// Before every instruction (and while one waits): a pending stop, then the fetch.
if Mode == User and IPEND == 1 then IPEND = 0; Trap(stop, 0); end
if PC MOD 4 != 0 or UInt(PC) >= ILS_SIZE then Trap(ifault, PC); end
Stall while PC overlaps the destination of an in-flight tgetx;
instr = ILS.Read(PC, 4);
\end{pseudo}

\section{Pseudo-instructions}

The assembler expands the pseudo-instructions of Table~\ref{tab:pseudo} into real instructions.

\begin{table}[htbp]
\centering\small
\caption{Assembler pseudo-instructions}\label{tab:pseudo}
\begin{tabularx}{\linewidth}{@{}>{\ttfamily}l >{\ttfamily}X@{}}
\toprule
\textrm{\bfseries Pseudo-instruction} & \textrm{\bfseries Expansion}\\\midrule
nop & addi x0, x0, 0\\
mv rd, rs & addi rd, rs, 0\\
neg rd, rs & sub rd, x0, rs\\
li rd, value & \textrm{1 instruction (}addi\textrm{) if it fits 16 signed bits; 2 (}lui; ori\textrm{) if it fits 32; else 6 (}lui; ori; slli 16; ori; slli 16; ori\textrm{)}\\
la rd, symbol & lui rd, hi; ori rd, rd, lo \textrm{(always 2 instructions; 0 $\le$ value $<2^{31}$)}\\
j label & jal x0, label\\
call label & jal ra, label\\
ret & jalr x0, ra, 0\\
jr rs & jalr x0, rs, 0\\
tfence & twait 0xffff\\
beqz / bnez / bltz / bgez rs, l & beq / bne / blt / bge rs, x0, l\\
bgtz rs, l \textrm{;} blez rs, l & blt x0, rs, l \textrm{;} bge x0, rs, l\\
bgt / ble rs, rt, l & blt / bge rt, rs, l\\
bgtu / bleu rs, rt, l & bltu / bgeu rt, rs, l\\
\bottomrule
\end{tabularx}
\end{table}
